\documentclass[10pt,a4paper]{amsart}
\usepackage[margin=19mm]{geometry}
\usepackage{amsmath,amssymb,mathtools}
\usepackage[scr=rsfs,cal=euler]{mathalfa}
\usepackage{xfrac}
\usepackage[unicode,hidelinks,bookmarks=false]{hyperref}
\newcommand{\ie}{i.e.\ }
\newcommand{\cf}{cf.\ }
\newcommand{\resp}{respectively\ }
\newcommand{\Subsection}{\subsection}
\providecommand{\nobreakdash}{\nobreak\hbox{-}}
\setlength{\emergencystretch}{2em}
\multlinegap0pt
\usepackage{amssymb}
\usepackage{enumerate}
\usepackage{xcolor}
\newtheorem{lemma}{Lemma}
\title{Exact characterizations for quantum conditional mutual information and some other entropies}
\author{Zhou Gang}
\date{Source: arXiv:2603.14650v2 (CC BY 4.0)}
\begin{document}
\maketitle

\noindent\emph{Source and licence.} Exact characterizations for quantum conditional mutual information and some other entropies. Version arXiv:2603.14650v2 (CC BY 4.0), distributed by arXiv under the Creative Commons Attribution 4.0 International licence, http://creativecommons.org/licenses/by/4.0/, as recorded in the arXiv licence metadata for this exact version and shown on the abstract page https://arxiv.org/abs/2603.14650v2. Full licence notice: \texttt{licenses/arxiv-2603.14650-CC-BY-4.0.txt}.\par\smallskip

\noindent\textit{Editorial scope.} Counted unit: the article's lemma eqn:DXOmega (source0.tex lines 1553--1562). The article's complete original proof of it is delivered with it (source0.tex lines 1563--1625, one \texttt{proof} environment of 63 lines). No mathematics is written, repaired or re-derived: the statement and the proof are the article's own lines, in the article's own notation, and no retained line is substituted or re-flowed.  No further source-local context is needed: every cross-reference of every retained line resolves inside the two delivered spans.  Nothing was omitted from any retained span, so the package carries no omission record.  No retained line contains a citation, so the wrapper carries no bibliography and no bibliography entry.  No retained line contains a figure environment, an image-inclusion command or drawing code, so no image is delivered and the package declares no asset dependency.  Editorial formatting only, in addition: an AMS \texttt{amsart} preamble that restates the article's own declarations and macros used by the retained lines, namely the environments \texttt{lemma} on separate counters, together with the standard packages those lines need.  The retained environments print in this document as Lemma 1 in order of appearance, whereas the article prints the same items with the numbering of its own full document.  The complete native source of the article as distributed in the arXiv e-print for this exact version is bundled unchanged as \texttt{sources/196423/source0.tex}.
\begin{lemma} 
Let $\Omega$ be a $N\times N$ Hermitian matrix, and $Y$ be a positive-definite $N\times N$ matrix. Then there exists a unique $N\times N$ Hermitian matrix $X$ such that
\begin{align}\label{eqn:DXOmega}
YX + XY = \Omega, 
\end{align} and it takes the form
\begin{equation}\label{eq:DXomega}
X = \int_{0}^{\infty} e^{-tY} \Omega e^{-tY} dt.
\end{equation}
In particular, if $\Omega > 0$ $(\geq 0)$ then $X > 0$ $(\geq 0)$.
\end{lemma}

\begin{proof}
Assuming the identity \eqref{eq:DXomega} holds and $\Omega\geq 0$, then since $e^{-tY} \Omega e^{-tY}\geq 0$ is for any $t$,
\[ X \geq 0. \]
By the same reason, if $\Omega > 0,\ (\geq 0)$, then $X > 0,\ (\geq 0)$.


What is left is to prove the identity \eqref{eq:DXomega}.

We start with considering an easy case, specifically, when $Y>0$ is a diagonal matrix, and hence, is of the form, for some $d_k>0$, $k=1,\cdots, N$,
\begin{equation}
  Y=\operatorname{diag}[d_1,d_2,\dots,d_N].
  \label{eqn:a}
\end{equation}

Denote the matrices $\Omega$ and $X$ by
\[
  \Omega=\bigl[w_{k,l}\bigr]_{k,l=1}^N,\ \text{and}\ 
  X=\bigl[x_{k,l}\bigr]_{k,l=1}^N.
\]
Then the equation takes the new form.
\[
  x_{k,l}\,(d_k+d_l)=w_{k,l}.
\]
Solve for $x_{k,l}$ to find
\[
  x_{k,l}
  =\frac{1}{d_k+d_l}\,w_{k,l}
  =\left(\int_0^\infty e^{-(d_k+d_l)t}\,dt\right) w_{k,l}.
\]
In the matrix form,
\begin{equation}\label{eqn:b}
  X=\int_0^\infty e^{-Yt}\,\Omega\,e^{-Yt}\,dt.
\end{equation}



Next, we consider the general case, i.e.\ $Y>0$, by converting it to the easier case \eqref{eqn:a} above.
Since $Y>0$, there exists a unitary matrix $U$ such that
\[
  \widetilde Y = UYU^*
\]
is diagonal, and hence is of the form \eqref{eqn:a}.

Apply $U$ and $U^*$ onto the appropriate places of equation \eqref{eqn:DXOmega} to find that
\[
  \widetilde Y\,\widetilde X+\widetilde X\,\widetilde Y=\widetilde\Omega,
\]
where $\widetilde X$ and $\widetilde\Omega$ are matrices defined as
\begin{equation}\label{eqn:XOmega}
  \widetilde X := UXU^*,
  \qquad
  \widetilde\Omega := U\Omega U^*.
\end{equation}

Since $\widetilde{D}$ is diagonal, the analysis for the easier case becomes applicable. By \eqref{eqn:b}, we find
\[
  \widetilde X = \int_0^\infty e^{-t\widetilde Y}\,\widetilde\Omega\,e^{-t\widetilde Y}\,dt.
\]
This and \eqref{eqn:XOmega}, together with applying $U^*$ and $U$ on the left and right, yield the desired identity
\[
  X = \int_0^\infty e^{-tY}\,\Omega\,e^{-tY}\,dt.
\]
\end{proof}

\end{document}
